\documentclass[12pt]{article}
\usepackage[a4paper, margin=2.5cm]{geometry}
\usepackage{graphicx}
\usepackage{amsmath}
\usepackage{amssymb}
\usepackage{amsthm}
\usepackage[T1]{fontenc}
\usepackage[utf8]{inputenc}
\usepackage{lmodern}
\usepackage{titlesec}
\usepackage{hyperref}
\usepackage{xurl}
\usepackage{tocloft}
\usepackage[spanish]{babel}
\usepackage{setspace}
\usepackage{array}
\usepackage{booktabs}
\usepackage{tikz}
\usetikzlibrary{automata, positioning, arrows.meta}
\setlength{\parindent}{0pt}
\hypersetup{hidelinks}
\onehalfspacing

\titleformat{\section}
  {\fontsize{13}{15}\selectfont\bfseries}
  {\thesection.}{0.5em}{}

\newtheorem{definicion}{Definición}[section]
\newtheorem{teorema}{Teorema}[section]
\newtheorem{proposicion}{Proposición}[section]
\newtheorem{ejemplo}{Ejemplo}[section]

\begin{document}
\thispagestyle{empty}
\begin{center}
    {\fontsize{14}{17}\selectfont\textbf{UNIVERSIDAD NACIONAL DE INGENIERÍA}}\\[0.2cm]
    {\fontsize{12}{15}\selectfont\textbf{FACULTAD DE CIENCIAS}}\\[0.2cm]
    {\fontsize{12}{15}\selectfont\textbf{ÁLGEBRA LINEAL I}}\\[0.4cm]
    {\includegraphics[width=8cm, height=8cm, keepaspectratio]{image2.png}}\\[0.5cm]
    {\fontsize{13}{16}\selectfont\textbf{PROYECTO 4 | AVANCE 3}}\\[0.3cm]
    {\fontsize{13}{16}\selectfont\textbf{ANÁLISIS Y PREDICCIÓN EN REDES MEDIANTE CADENAS DE MARKOV Y
            ÁLGEBRA LINEAL}}\\[0.2cm]
    {\fontsize{11}{14}\selectfont\textit{Aplicación a la movilidad de usuarios en el transporte público
            de Lima Metropolitana}}\\[0.2cm]
    {\fontsize{11}{14}\selectfont\textbf{Integrantes del Grupo 1}}\\[0.3cm]
    \begin{tabular}{c p{12cm} c}
        \textbf{N°} & \textbf{Apellidos y Nombres}          & \textbf{Código} \\[0.3cm]
        1           & ALTEZ INFANTAS, JUAN CARLOS           & 20261602I       \\[0.3cm]
        2           & BARDALES SUCLUPE, WILY JUNIOR FILOMON & 20241481A       \\[0.3cm]
        3           & GARCIA SAMAN, ARON PATRICK            & 20252147K       \\[0.3cm]
        4           & HUALLPA TORRES, YAN CESAR             & 20250554H       \\[0.3cm]
        5           & POMASUNCO HUILLCA, RODRIGO ANDRE      & 20251442I       \\[0.3cm]
        6           & TIRADO NINA, HECTOR EDUARDO           & 20250475K       \\[0.3cm]
    \end{tabular}\vspace{0.2cm}

    {\fontsize{11}{14}\selectfont\textbf{Docente: PÉREZ CUPE, RÓSULO HILARIÓN}}\\[0.2cm]
    {\fontsize{11}{14}\selectfont\textbf{2026 - II}}\\[0.2cm]
    {\fontsize{11}{14}\selectfont\textbf{LIMA -- PERÚ}}\\[0.2cm]
\end{center}

\newpage
\tableofcontents
\newpage

\section*{RESUMEN}
Se modela la movilidad de usuarios en el transporte público de Lima Metropolitana mediante una
cadena de Markov en tiempo discreto. Las zonas de la ciudad son los estados y los desplazamientos
entre ellas, las transiciones. La dinámica se representa con una matriz estocástica por filas
y vectores de estado. En este avance se precisan los estados, las variables y los supuestos del
modelo, se explica cómo se estimará la matriz de transición a partir de datos institucionales y
se desarrolla la teoría de Álgebra Lineal necesaria (potencias de matrices, sistemas lineales,
valores y vectores propios, teorema de Perron--Frobenius) para analizar la evolución del sistema y
su comportamiento de largo plazo \cite{norris,meyer}.

\textbf{Palabras clave:} cadena de Markov, matriz estocástica, distribución estacionaria,
valores propios, movilidad urbana, álgebra lineal.

\section{INTRODUCCIÓN}
Las cadenas de Markov permiten representar sistemas que evolucionan entre estados mediante
probabilidades de transición \cite{norris,kemeny}. En este proyecto se aplica ese enfoque a la
movilidad de usuarios entre zonas de Lima Metropolitana. El interés para Álgebra Lineal radica
en que toda la dinámica puede expresarse con vectores y matrices: las potencias de la matriz de
transición describen la evolución temporal y los valores propios permiten estudiar el
comportamiento de largo plazo \cite{meyer,strang}.

Este documento corresponde al segundo avance. Respecto al primero, se concretan los estados y
los supuestos, se desarrolla el marco teórico y se muestra, con un ejemplo numérico, el análisis
que se realizará sobre la matriz de transición del modelo. La implementación computacional se
presentará en el tercer y último avance.

\section{DESCRIPCIÓN DEL PROBLEMA}
\subsection{Contexto}
Lima Metropolitana concentra una gran parte de la población del país y depende en buena medida
del transporte público para sus desplazamientos diarios. Este sistema reúne varios modos
(Metro de Lima, Metropolitano, corredores complementarios y buses), cuya regulación en Lima y
Callao corresponde a la Autoridad de Transporte Urbano (ATU) \cite{atu}. Conocer cómo se
redistribuyen los usuarios entre zonas es útil para anticipar la congestión, dimensionar la
oferta de rutas y evaluar la planificación del sistema.

\subsection{Planteamiento}
Se selecciona un conjunto finito de zonas y cada desplazamiento de un usuario se considera una
transición entre estados. La pregunta central es:

\vspace{0.3cm}
{\itshape ¿Cómo evoluciona la distribución de usuarios entre distintas zonas de Lima
    Metropolitana cuando los desplazamientos se modelan mediante una cadena de Markov, y qué
    distribución se alcanza a largo plazo?}
\vspace{0.3cm}

\subsection{Alcance}
El modelo trabaja a nivel de \textbf{zonas} (y no de estaciones o rutas individuales), porque ese nivel de agregación mantiene el modelo manejable y es el habitual en los estudios de movilidad urbana. Un \textbf{paso} $k \to k+1$ corresponde a \textbf{un viaje}: $X_k$ es la zona donde se encuentra el usuario y la transición describe hacia qué zona se dirige en su siguiente viaje.

\newpage

\section{IMPORTANCIA DEL ÁLGEBRA LINEAL EN EL PROBLEMA}
En una cadena de Markov finita, las probabilidades de transición se organizan en una matriz y la
distribución del sistema en un vector \cite{meyer,strang}. Cada herramienta del curso resuelve una
parte concreta del problema:
\begin{itemize}
    \item \textbf{Producto matriz--vector:} calcula la distribución del siguiente periodo,
          $x^{(k+1)}=x^{(k)}P$.
    \item \textbf{Potencias de matrices:} permiten predecir a $k$ pasos, $x^{(k)}=x^{(0)}P^k$.
    \item \textbf{Sistemas lineales homogéneos:} la distribución estacionaria se obtiene
          resolviendo $\pi(P-I)=0$ con $\sum_i\pi_i=1$.
    \item \textbf{Valores y vectores propios:} $\pi$ es un vector propio izquierdo asociado a
          $\lambda=1$, y el resto del espectro controla la velocidad de convergencia.
    \item \textbf{Diagonalización:} si $P=VDV^{-1}$, entonces $P^k=VD^kV^{-1}$, lo que
          simplifica el cálculo de potencias \cite{lay,meyer}.
\end{itemize}
Por ello el Álgebra Lineal constituye el núcleo matemático del modelo.

\section{JUSTIFICACIÓN DE LA ELECCIÓN DEL PROBLEMA}
El problema fue elegido porque:
\begin{itemize}
    \item admite una representación clara mediante estados y transiciones;
    \item se expresa directamente con vectores y matrices;
    \item permite aplicar contenidos centrales del curso (sistemas lineales, valores propios,
          diagonalización);
    \item dispone de fuentes institucionales para obtener datos \cite{atu,inei}.
\end{itemize}
Además, las cadenas de Markov son una herramienta estándar para modelar sistemas probabilísticos
discretos \cite{norris,ross}.
\newpage

\section{OBJETIVOS}
\subsection{Objetivo general}
Modelar la movilidad de usuarios entre zonas de Lima Metropolitana mediante una cadena de Markov
finita y analizar su evolución y comportamiento de largo plazo con herramientas de Álgebra Lineal.

\subsection{Objetivos específicos}
\begin{itemize}
    \item Definir los estados, las variables y los supuestos del modelo.
    \item Identificar y documentar los datos de movilidad necesarios para estimar las transiciones.
    \item Construir la matriz de transición estocástica $P$ a partir de frecuencias de viaje.
    \item Representar la distribución de usuarios mediante vectores de estado.
    \item Verificar la regularidad de $P$ y justificar la existencia y unicidad de la distribución
          estacionaria.
    \item Calcular la distribución estacionaria y los valores propios de $P$ e interpretarlos en el
          contexto de Lima.
\end{itemize}

\newpage
\section{ESTADOS, VARIABLES, DATOS Y SUPUESTOS}

\subsection{Estados}
Se considera un conjunto finito $S=\{1,2,\ldots,n\}$ donde cada estado es una zona de Lima
Metropolitana. Como esquema de trabajo se consideran cinco grandes zonas:
\[
    1:\ \text{Lima Norte},\quad 2:\ \text{Lima Centro},\quad 3:\ \text{Lima Este},\quad
    4:\ \text{Lima Sur},\quad 5:\ \text{Callao}.
\]
Cada zona agrupa un conjunto de distritos contiguos con características de movilidad similares. La selección definitiva dependerá de la disponibilidad y del nivel de detalle de los datos de movilidad.

\subsection{Variables}
\begin{itemize}
    \item $X_k$: zona ocupada por un usuario en el periodo $k$;
    \item $n_{ij}$: número de desplazamientos observados de la zona $i$ a la zona $j$;
    \item $p_{ij}$: probabilidad de transición de $i$ a $j$;
    \item $P=(p_{ij})\in\mathbb{R}^{n\times n}$: matriz de transición;
    \item $x^{(k)}=(x^{(k)}_1,\ldots,x^{(k)}_n)$: distribución de usuarios entre zonas en el
          periodo $k$ (vector fila);
    \item $\pi$: distribución estacionaria.
\end{itemize}

\subsection{Datos y estimación de la matriz}
Las entradas de $P$ se estiman a partir de frecuencias de desplazamiento entre zonas (matriz
origen--destino) obtenidas de reportes de ATU e INEI \cite{atu,inei}:
\[
    p_{ij}=\frac{n_{ij}}{\sum_{j=1}^{n} n_{ij}},\qquad i,j=1,\ldots,n.
\]
Con esta definición cada fila suma 1 por construcción, de modo que $P$ es estocástica por filas.

Los conteos $n_{ij}$ provendrán de matrices origen--destino o frecuencias de desplazamiento entre zonas publicadas por ATU e INEI. Mientras tanto, el procedimiento de análisis se desarrolla en la sección~\ref{sec:ejemplo} con una matriz ilustrativa, que sigue exactamente los mismos pasos que se aplicarán a la matriz estimada.

\subsection{Supuestos del modelo}
Se declaran explícitamente los siguientes supuestos, cuya validez se discute en las conclusiones:
\begin{enumerate}
    \item \textbf{Propiedad de Markov:} la zona de destino depende solo de la zona actual y no del
          historial previo de viajes.
    \item \textbf{Homogeneidad en el tiempo:} la matriz $P$ es la misma en todos los periodos
          (no se distingue entre hora punta y hora valle).
    \item \textbf{Población cerrada:} los usuarios del sistema no entran ni salen de la red
          considerada; la suma de $x^{(k)}$ se mantiene igual a 1.
    \item \textbf{Viajes internos:} los desplazamientos dentro de una misma zona se contabilizan
          como transiciones $i\to i$, de modo que $p_{ii}$ puede ser positivo.
    \item \textbf{Representatividad de los datos:} las frecuencias institucionales representan el
          comportamiento general de los usuarios.
\end{enumerate}

\newpage
\section{PRELIMINARES MATEMÁTICOS}

\subsection{Cadena de Markov}
Una cadena de Markov en tiempo discreto es un proceso $(X_k)_{k\ge 0}$ con valores en $S$ que
satisface
\[
    \mathbb{P}(X_{k+1}=j \mid X_k=i, X_{k-1},\ldots,X_0)=\mathbb{P}(X_{k+1}=j \mid X_k=i),
\]
es decir, el siguiente estado depende solo del estado actual \cite{norris,kemeny}.

\subsection{Matriz de transición}
\begin{definicion}
    Una matriz $P=(p_{ij})\in\mathbb{R}^{n\times n}$ es \emph{estocástica por filas} si
    $p_{ij}\ge 0$ para todo $i,j$ y $\sum_{j=1}^{n}p_{ij}=1$ para todo $i$.
\end{definicion}
En el modelo, $p_{ij}=\mathbb{P}(X_{k+1}=j\mid X_k=i)$ \cite{norris,ross}. Equivalentemente,
$P\mathbf{1}=\mathbf{1}$, donde $\mathbf{1}$ es el vector columna de unos.

\begin{proposicion}\label{prop:producto}
    Si $P$ y $Q$ son estocásticas por filas de orden $n$, entonces $PQ$ también lo es. En
    particular, $P^k$ es estocástica para todo $k\ge 1$.
\end{proposicion}
\begin{proof}
    Las entradas de $PQ$ son sumas de productos de números no negativos, luego son no negativas.
    Además, $(PQ)\mathbf{1}=P(Q\mathbf{1})=P\mathbf{1}=\mathbf{1}$.
\end{proof}

\subsection{Vector de estado y evolución}
El vector fila $x^{(k)}=(x^{(k)}_1,\ldots,x^{(k)}_n)$, con $x^{(k)}_i\ge 0$ y
$\sum_i x^{(k)}_i=1$, representa la distribución de usuarios en el periodo $k$. Su evolución es
\[
    x^{(k+1)}=x^{(k)}P,\qquad\text{y por inducción}\qquad x^{(k)}=x^{(0)}P^k.
\]
Además, $p^{(k)}_{ij}=(P^k)_{ij}$ es la probabilidad de pasar de $i$ a $j$ en exactamente $k$
pasos \cite{meyer,strang}.

\subsection{Clasificación de estados}
\begin{definicion}
    El estado $j$ es \emph{accesible} desde $i$ si existe $k\ge 0$ con $(P^k)_{ij}>0$. Los
    estados $i$ y $j$ se \emph{comunican} si cada uno es accesible desde el otro. La cadena es
    \emph{irreducible} si todos los estados se comunican entre sí.
\end{definicion}
\begin{definicion}
    El \emph{periodo} de un estado $i$ es $d_i=\operatorname{mcd}\{k\ge 1: (P^k)_{ii}>0\}$. Si
    $d_i=1$ el estado es \emph{aperiódico}; en una cadena irreducible todos los estados tienen el
    mismo periodo.
\end{definicion}
\begin{definicion}
    Una matriz estocástica $P$ es \emph{regular} si existe $m\ge 1$ tal que todas las entradas de
    $P^m$ son estrictamente positivas.
\end{definicion}
Una cadena finita es regular si y solo si es irreducible y aperiódica \cite{seneta,levin}. Un criterio
práctico: si $P$ es irreducible y algún $p_{ii}>0$, entonces es regular.

\subsection{Distribución estacionaria}
\begin{definicion}
    Una distribución $\pi$ es \emph{estacionaria} si
    \[
        \pi P=\pi,\qquad \pi_i\ge 0,\qquad \sum_{i=1}^n\pi_i=1.
    \]
\end{definicion}
La condición equivale al sistema lineal homogéneo $\pi(P-I)=0$, es decir, $\pi$ es un
\textbf{vector propio izquierdo} de $P$ asociado a $\lambda=1$ (o, de forma equivalente, un vector
propio de $P^{\top}$) \cite{norris,seneta}.

\subsection{Valores propios y convergencia}
\begin{proposicion}\label{prop:espectro}
    Si $P$ es estocástica por filas, entonces $\lambda=1$ es valor propio de $P$ (con vector propio
    derecho $\mathbf{1}$) y todo valor propio $\lambda$ cumple $|\lambda|\le 1$.
\end{proposicion}
\begin{proof}
    Como $P\mathbf{1}=\mathbf{1}$, se tiene $\lambda=1$. Sea $Pv=\lambda v$ con $v\ne 0$ y sea $i$
    un índice con $|v_i|$ máximo. Entonces
    $|\lambda||v_i|=\bigl|\sum_j p_{ij}v_j\bigr|\le\sum_j p_{ij}|v_j|\le|v_i|$, de donde
    $|\lambda|\le 1$.
\end{proof}

\begin{teorema}[Perron--Frobenius para matrices estocásticas regulares]\label{teo:pf}
    Sea $P$ estocástica por filas y regular. Entonces:
    \begin{enumerate}
        \item $\lambda=1$ es un valor propio simple y todo otro valor propio cumple $|\lambda|<1$;
        \item existe una única distribución estacionaria $\pi$ y todas sus entradas son positivas;
        \item $P^k\to\mathbf{1}\pi$ cuando $k\to\infty$; en consecuencia, $x^{(k)}\to\pi$ para
              cualquier distribución inicial $x^{(0)}$.
    \end{enumerate}
\end{teorema}
La demostración puede consultarse en \cite{seneta,levin}. La velocidad de la convergencia está
gobernada por $|\lambda_2|$, el segundo mayor módulo entre los valores propios: cuanto menor es
$|\lambda_2|$, más rápido $x^{(k)}$ se aproxima a $\pi$.

\subsection{Diagonalización}
Si $P$ es diagonalizable, $P=VDV^{-1}$ con $D=\operatorname{diag}(1,\lambda_2,\ldots,\lambda_n)$,
entonces $P^k=VD^kV^{-1}$ y
\[
    x^{(k)}=x^{(0)}VD^kV^{-1}.
\]
Como $|\lambda_i|<1$ para $i\ge 2$, $D^k\to\operatorname{diag}(1,0,\ldots,0)$, lo que explica
el teorema~\ref{teo:pf} \cite{meyer,lay}.

\newpage
\section{EJEMPLO ILUSTRATIVO DEL ANÁLISIS}\label{sec:ejemplo}
Para mostrar el procedimiento se usa una matriz \textbf{ilustrativa} de 3 zonas, con valores elegidos solo para explicar el método. En el siguiente avance se reemplazará por la matriz estimada a partir de los datos de movilidad, siguiendo exactamente el mismo procedimiento.

Sean las zonas $A$, $B$, $C$ y
\[
    P=\begin{pmatrix}
        0{,}7 & 0{,}2 & 0{,}1 \\
        0{,}3 & 0{,}5 & 0{,}2 \\
        0{,}2 & 0{,}3 & 0{,}5
    \end{pmatrix},\qquad x^{(0)}=(0{,}5\;\;0{,}3\;\;0{,}2).
\]

\begin{figure}[h]
    \centering
    \begingroup
    \shorthandoff{>}
    \begin{tikzpicture}[->, >=Stealth, node distance=3.2cm, auto, thick,
            every state/.style={minimum size=1cm}]
        \node[state] (A) {$A$};
        \node[state] (B) [right of=A] {$B$};
        \node[state] (C) [below right=2.2cm and 1.6cm of A] {$C$};
        \path
        (A) edge[loop above] node {0,7} (A)
        (B) edge[loop above] node {0,5} (B)
        (C) edge[loop below] node {0,5} (C)
        (A) edge[bend left=15] node {0,2} (B)
        (B) edge[bend left=15] node {0,3} (A)
        (A) edge[bend left=10] node[pos=0.6,left] {0,1} (C)
        (C) edge[bend left=10] node[pos=0.4,right] {0,2} (A)
        (B) edge[bend left=10] node[pos=0.6,right] {0,2} (C)
        (C) edge[bend left=10] node[pos=0.4,below] {0,3} (B);
    \end{tikzpicture}
    \endgroup
    \caption{Grafo de transiciones del ejemplo ilustrativo.}
\end{figure}

\subsection{Verificación de \texorpdfstring{$P$}{P}}
Todas las entradas son no negativas y cada fila suma 1, luego $P$ es estocástica. Como todos los
$p_{ij}>0$, la cadena es irreducible y aperiódica; es decir, $P$ es regular y se aplica el
teorema~\ref{teo:pf}.

\subsection{Evolución de la distribución}
\begin{align*}
    x^{(1)} & =x^{(0)}P=(0{,}480\;\;0{,}310\;\;0{,}210), \\
    x^{(2)} & =x^{(1)}P=(0{,}471\;\;0{,}314\;\;0{,}215).
\end{align*}

\subsection{Distribución estacionaria}
Se resuelve $\pi P=\pi$ con $\sum_i\pi_i=1$. Las ecuaciones independientes son
\begin{align*}
    0{,}3\pi_1 & =0{,}3\pi_2+0{,}2\pi_3, \\
    0{,}5\pi_3 & =0{,}1\pi_1+0{,}2\pi_2.
\end{align*}
De la segunda, $\pi_3=0{,}2\pi_1+0{,}4\pi_2$. Al sustituir en la primera se obtiene
$0{,}26\pi_1=0{,}38\pi_2$, es decir, $\pi_2=\tfrac{13}{19}\pi_1$ y $\pi_3=\tfrac{9}{19}\pi_1$.
Con la condición de normalización:
\[
    \pi=\left(\tfrac{19}{41},\;\tfrac{13}{41},\;\tfrac{9}{41}\right)\approx
    (0{,}4634\;\;0{,}3171\;\;0{,}2195).
\]
El mismo resultado se obtiene por eliminación de Gauss--Jordan. Multiplicando por 10 las dos primeras ecuaciones y reemplazando la tercera (redundante) por la condición de normalización, la matriz aumentada se reduce así:
\[
    \left(\begin{array}{ccc|c}
        -3 & 3  & 2 & 0 \\
        2  & -5 & 3 & 0 \\
        1  & 1  & 1 & 1
    \end{array}\right)
    \;\longrightarrow\;
    \left(\begin{array}{ccc|c}
        1 & 0 & 0 & 19/41 \\
        0 & 1 & 0 & 13/41 \\
        0 & 0 & 1 & 9/41
    \end{array}\right).
\]
Nótese que $x^{(2)}$ ya está muy cerca de $\pi$. El tiempo medio de retorno a la zona $i$ es
$1/\pi_i$; por ejemplo, para $A$ es $41/19\approx 2{,}16$ pasos.

\subsection{Valores propios}
La traza es $1{,}7$ y el determinante es $0{,}11$. Como uno de los valores propios es $1$, los
otros dos satisfacen $\lambda^2-0{,}7\lambda+0{,}11=0$, de donde
\[
    \lambda_1=1,\qquad \lambda_{2,3}=\frac{0{,}7\pm\sqrt{0{,}05}}{2}\approx 0{,}462,\;0{,}238.
\]
Como $|\lambda_2|\approx 0{,}46<1$, la convergencia a $\pi$ es rápida: el error se reduce
aproximadamente a la mitad en cada paso.
\newpage
\section{BIBLIOGRAFÍA}
\begin{thebibliography}{10}
    \bibitem{norris} Norris, J. R. (1998). \emph{Markov Chains}. Cambridge University Press. \url{https://doi.org/10.1017/CBO9780511819181}
    \bibitem{meyer} Meyer, C. D. (2023). \emph{Matrix Analysis and Applied Linear Algebra}. 2nd ed., SIAM. \url{https://doi.org/10.1137/1.9781611972082}
    \bibitem{ross} Ross, S. M. (2019). \emph{Introduction to Probability Models}. 12th ed., Academic. \url{https://www.sciencedirect.com/book/9780128143469/introduction-to-probability-models}
    \bibitem{kemeny} Kemeny, J. G., \& Snell, J. L. (1960). \emph{Finite Markov Chains}. Van Nostrand.
    \bibitem{seneta} Seneta, E. (2006). \emph{Non-negative Matrices and Markov Chains}. 3rd ed., Springer.
    \bibitem{levin} Levin, D. A., Peres, Y., \& Wilmer, E. L. (2009). \emph{Markov Chains and Mixing Times}. American Mathematical Society.
    \bibitem{strang} Strang, G. (2016). \emph{Introduction to Linear Algebra}. 5th ed., Wellesley--Cambridge Press.
    \bibitem{lay} Lay, D. C., Lay, S. R., \& McDonald, J. J. (2016). \emph{Linear Algebra and Its Applications}. 5th ed., Pearson.
    \bibitem{atu} Autoridad de Transporte Urbano para Lima y Callao (ATU). (2023). \emph{Reportes de demanda y estadísticas de transporte}. \url{https://www.atu.gob.pe}
    \bibitem{inei} Instituto Nacional de Estadística e Informática (INEI). (2023). \emph{Estadísticas de movilidad y población}. \url{https://www.inei.gob.pe}
\end{thebibliography}
\end{document}